% Original schematic: three DAGs and their common CPDAG.
\begin{tikzpicture}[x=1mm,y=1mm,font=\small,
  pgm variable/.style={circle,draw=FigureInk,fill=FigurePaper,
    minimum size=7mm,inner sep=0pt},
  edge/.style={-{Stealth[length=2mm]},draw=FigureInk,thick},
  heading/.style={font=\figurefont\footnotesize,text=FigureInk}]
  \node[heading] at (54,76) {正向链};
  \node[pgm variable] (a1) at (18,65) {$X_1$};
  \node[pgm variable] (a2) at (54,65) {$X_2$};
  \node[pgm variable] (a3) at (90,65) {$X_3$};
  \draw[edge] (a1)--(a2);
  \draw[edge] (a2)--(a3);

  \node[heading] at (54,54) {分叉};
  \node[pgm variable] (b1) at (18,43) {$X_1$};
  \node[pgm variable] (b2) at (54,43) {$X_2$};
  \node[pgm variable] (b3) at (90,43) {$X_3$};
  \draw[edge] (b2)--(b1);
  \draw[edge] (b2)--(b3);

  \node[heading] at (54,32) {反向链};
  \node[pgm variable] (c1) at (18,21) {$X_1$};
  \node[pgm variable] (c2) at (54,21) {$X_2$};
  \node[pgm variable] (c3) at (90,21) {$X_3$};
  \draw[edge] (c2)--(c1);
  \draw[edge] (c3)--(c2);

  \draw[FigureGrid] (8,11)--(100,11);
  \node[heading,text=FigureYellow] at (54,3) {共同的CPDAG};
  \node[pgm variable,draw=FigureYellow] (p1) at (18,-9) {$X_1$};
  \node[pgm variable,draw=FigureYellow] (p2) at (54,-9) {$X_2$};
  \node[pgm variable,draw=FigureYellow] (p3) at (90,-9) {$X_3$};
  \draw[FigureYellow,thick] (p1)--(p2)--(p3);
  \node[font=\footnotesize,text=FigureMuted] at (54,-21)
    {$X_1\perp X_3\mid X_2$；不确定方向并非独立选择};
\end{tikzpicture}
